% IN5340/IN9340 Project IV Presentation
% Build (from this folder): pdflatex -interaction=nonstopmode -halt-on-error presentation_p4.tex

\documentclass[aspectratio=169,10pt]{beamer}

\usetheme{Madrid}
\usecolortheme{default}

\usepackage{graphicx}
\usepackage{amsmath,amssymb}
\usepackage{booktabs}
\usepackage{mathtools}
\usepackage[T1]{fontenc}
\usepackage[utf8]{inputenc}
\usepackage{hyperref}

\graphicspath{{figs/}}

\newcommand{\fullfig}[2][]{%
\centering
\includegraphics[width=\linewidth,height=0.78\textheight,keepaspectratio,#1]{#2}%
}
\newcommand{\colfig}[2][]{%
\includegraphics[width=\linewidth,height=0.58\textheight,keepaspectratio,#1]{#2}%
}

\title[IN5340 Project IV]{IN5340/IN9340 Project IV: Spectral Estimation}
\author{Theodor W\aa lberg}
\institute{University of Oslo}
\date{March 2026}

\begin{document}

\begin{frame}
\titlepage
\vspace{-1.0em}
\centering
\footnotesize
Code: \href{https://github.com/Twaal/Statistical_Signal_Processing_IN5340/tree/master/Project_4}{\nolinkurl{github.com/Twaal/Statistical_Signal_Processing_IN5340/tree/master/Project_4}}
\end{frame}

\begin{frame}{Agenda}
\tableofcontents
\end{frame}

% =============================================================================
\section{Introduction}
% =============================================================================

\begin{frame}{Problem statement}
\begin{itemize}
\item Sonar data from HUGIN AUV: $N_h=32$ hydrophones, single ping.
\item Data matrix $\mathbf{X}\in\mathbb{C}^{N_t\times 32}$, sampling frequency $f_s=40\,$kHz.
\item Transmitted signal: LFM pulse, bandwidth $B=30\,$kHz, centre frequency $f_c=100\,$kHz.
\[
s_{Tx}(t) = \exp\!\bigl(j2\pi\alpha t^2/2\bigr),\quad |t|\le T_p/2, \qquad \alpha = B/T_p
\]
\item Received data are basebanded (carrier removed).
\item Goal: estimate and compare power spectral densities using periodogram, modified periodogram, Welch, and multitaper methods; analyse spectrograms.
\end{itemize}
\end{frame}

% =============================================================================
\section{Exercise 1A - Spectral Estimation (Channel 14)}
% =============================================================================

\begin{frame}[fragile]{1A: Periodogram and modified periodogram}
\scriptsize
\begin{verbatim}
# periodogram
fft_result = np.fft.fft(channel_14_data)
periodogram = (1/(N*fs)) * np.abs(fft_result)**2
periodogram_shifted = np.fft.fftshift(periodogram)

# modified periodogram with Hanning window
win = np.hanning(N)
U = np.mean(win**2)
windowed = channel_14_data * win
fft_win = np.fft.fft(windowed)
periodogram_windowed = (1/(N*fs*U)) * np.abs(fft_win)**2
\end{verbatim}
\end{frame}

\begin{frame}[fragile]{1A: Welch method and multitaper}
\scriptsize
\begin{verbatim}
# Welch (L=256, D=128, Hanning)
welch_psd = np.zeros(L)
for i in range(num_segments):
    start = i*(L - D)
    seg = channel_14_data[start:start+L] * welch_win
    welch_psd += (1/(L*fs*U_welch)) * np.abs(np.fft.fft(seg))**2
welch_psd /= num_segments

# Multitaper (NW=4, K=7, DPSS tapers)
tapers, eigenvalues = windows.dpss(N, NW, K, return_ratios=True)
multitaper_psd = np.zeros(N)
for k in range(K):
    multitaper_psd += (1/fs) * np.abs(np.fft.fft(channel_14_data * tapers[k]))**2
multitaper_psd /= K
\end{verbatim}
\end{frame}

\begin{frame}{1A: Comparison of spectral estimates}
\begin{columns}[T,onlytextwidth]
\column{0.42\textwidth}
\small
\begin{itemize}
\item Channel 14, $M=1200$, $N=1024$.
\item Four estimators plotted: periodogram, modified (Hanning), Welch ($L=256$, $D=128$), multitaper ($NW=4$, $K=7$).
\item The LFM passband ($\pm15\,$kHz) is visible in all estimates.
\item Periodogram shows the most variance; Welch and multitaper are smoother.
\end{itemize}
\column{0.58\textwidth}
\colfig{plot_06.png}
\end{columns}
\end{frame}

\begin{frame}{1A: Bias and variance analysis}
\small
\textbf{Does bias reduction help?}\\
The boxcar periodogram is theoretically unbiased only for infinite data. With finite $N=1024$, spectral leakage introduces bias. All methods show negative mean bias in the passband:
\vspace{0.3em}

\begin{tabular}{lr}
\toprule
Method & Mean Bias (dB) \\
\midrule
Periodogram & $-3.86$ \\
Modified (Hanning) & $-3.35$ \\
Welch ($L\!=\!256$) & $-1.57$ \\
Multitaper ($K\!=\!7$) & $-1.43$ \\
\bottomrule
\end{tabular}

\vspace{0.5em}
\textbf{Does variance reduction help?}\\
\begin{tabular}{lr}
\toprule
Method & Mean Variance \\
\midrule
Periodogram & $9.19$ \\
Modified (Hanning) & $11.17$ \\
Welch ($L\!=\!256$) & $0.81$ \\
Multitaper ($K\!=\!7$) & $0.75$ \\
\bottomrule
\end{tabular}

\vspace{0.3em}
Welch and multitaper reduce variance by $>10\times$, giving a much more reliable estimate.
\end{frame}

\begin{frame}{1A: Estimators vs ideal LFM spectrum}
\fullfig{plot_07.png}
\end{frame}

\begin{frame}{1A: Bandwidth estimation and spectrum issues}
\small
\textbf{Bandwidth estimation} (multitaper, $-6\,$dB from passband level):
\begin{itemize}
\item Passband level: $3.45\,$dB
\item $-6\,$dB threshold: $-2.55\,$dB
\item Estimated edges: $-17.58\,$kHz to $14.53\,$kHz
\item \textbf{Estimated BW: $32.11\,$kHz} (ideal: $30\,$kHz)
\end{itemize}

\vspace{0.5em}
\textbf{What is wrong with the spectrum?}
\begin{itemize}
\item Spectral leakage widens the apparent passband beyond the true $30\,$kHz.
\item An interfering tone near $-17\,$kHz occupies the lower band and skews the bandwidth estimate upward.
\end{itemize}
\end{frame}

% =============================================================================
\section{Exercise 1B - Spectral Estimation (Channel 9)}
% =============================================================================

\begin{frame}{1B: Comparison of spectral estimates}
\begin{columns}[T,onlytextwidth]
\column{0.40\textwidth}
\small
\begin{itemize}
\item Channel 9, $M=7000$, $N=2048$.
\item Signal dominated by additive noise and narrowband interference from "tonal lines".
\item Periodogram resolves closely spaced tones near $16$-$17\,$kHz that merge in Welch/multitaper.
\end{itemize}
\column{0.60\textwidth}
\colfig{plot_12.png}
\end{columns}
\end{frame}

\begin{frame}{1B: Discussion}
\small
\textbf{1. Does bias reduction help?}\\
No. The signal is dominated by narrowband tones in noise, not a wideband LFM pulse. There is no well-defined passband where spectral leakage matters, so bias reduction provides no practical benefit.

\vspace{0.5em}
\textbf{2. Does variance reduction help?}\\
Variance reduction is counterproductive here. At $16$-$17\,$kHz two closely spaced tones are distinguishable in the high-variance periodogram but merge in the smoothed Welch estimate.

\vspace{0.5em}
\textbf{3. Consequence of variance reduction?}\\
A bias-variance-resolution tradeoff: smoothing reduces variance but also reduces frequency resolution. The Welch method with $L=256$ gives $\Delta f = f_s/L \approx 156\,$Hz, which is too coarse to separate nearby tones.
\end{frame}

% =============================================================================
\section{Exercise 2A - Spectrogram Analysis (Channel 14)}
% =============================================================================

\begin{frame}{2A: Timeseries - Channel 14}
\begin{columns}[T,onlytextwidth]
\column{0.38\textwidth}
\small
\begin{itemize}
\item Channel 14, $M=400$, $N=8192$.
\item Several distinct events with varying amplitude and duration in additive noise.
\item Not WSS over the full window motivates time-frequency analysis.
\end{itemize}
\column{0.62\textwidth}
\colfig{plot_13.png}
\end{columns}
\end{frame}

\begin{frame}{2A: Welch PSD - Channel 14}
\begin{columns}[T,onlytextwidth]
\column{0.38\textwidth}
\small
\begin{itemize}
\item Welch with $L=512$, $D=256$, Hanning.
\item Several tonal clusters visible.
\item Original LFM passband present but tones dominate.
\item The global PSD hides the time-varying nature of these events.
\end{itemize}
\column{0.62\textwidth}
\colfig{plot_14.png}
\end{columns}
\end{frame}

\begin{frame}[fragile]{2A: STFT implementation}
\scriptsize
\begin{verbatim}
for idx, Ls in enumerate(seg_lengths):      # L = 64, 128, 256, 512
    Ds = Ls // 2
    Nfft = Ls * zp                          # zero-pad x4
    w = np.hanning(Ls);  Uw = np.mean(w**2)

    nseg = (N_2a - Ls) // Ds + 1
    stft = np.zeros((Nfft, nseg))
    for i in range(nseg):
        s = i * Ds
        seg = ch14_long[s:s+Ls] * w
        padded = np.zeros(Nfft, dtype=complex)
        padded[:Ls] = seg
        stft[:, i] = (1/(Ls*fs*Uw)) * np.abs(np.fft.fft(padded))**2

    stft_dB = 10*np.log10(np.fft.fftshift(stft, axes=0) + 1e-30)
\end{verbatim}
\end{frame}

\begin{frame}{2A: STFT - four segment lengths}
\fullfig{plot_15.png}
\end{frame}

\begin{frame}{2A: Discussion}
\small
\textbf{1. Consequence of segment length?}\\
Segment length controls the time-frequency resolution tradeoff:
\begin{itemize}
\item Short $L$ (64): good time resolution ($\Delta t = 1.6\,$ms), poor frequency resolution ($\Delta f = 625\,$Hz).
\item Long $L$ (512): good frequency resolution ($\Delta f = 78\,$Hz), poor time resolution ($\Delta t = 12.8\,$ms).
\end{itemize}

\vspace{0.5em}
\textbf{2. Interfering chips (100-200\,ms)?}\\
\begin{itemize}
\item The chips are approximately $20\,$ms in duration.
\item Their centre frequencies show a slight shift from chip to chip, possibly due to Doppler from a radially accelerating interferer.
\end{itemize}
\end{frame}

% =============================================================================
\section{Exercise 2B - Spectrogram Analysis (Channel 32)}
% =============================================================================

\begin{frame}{2B: Timeseries - Channel 32}
\begin{columns}[T,onlytextwidth]
\column{0.38\textwidth}
\small
\begin{itemize}
\item Channel 32, $M=400$, $N=8192$.
\item Real part of the signal plotted.
\item A burst of energy around $20\,$ms followed by decaying amplitude.
\end{itemize}
\column{0.62\textwidth}
\colfig{plot_16.png}
\end{columns}
\end{frame}

\begin{frame}{2B: Welch PSD - Channel 32}
\begin{columns}[T,onlytextwidth]
\column{0.38\textwidth}
\small
\begin{itemize}
\item Welch with $L=512$, $D=256$, Hanning.
\item Spectrum resembles the LFM passband ($\pm15\,$kHz) with some leakage.
\item Weak interfering tone near $-17\,$kHz.
\item Much cleaner than Channel 14.
\end{itemize}
\column{0.62\textwidth}
\colfig{plot_17.png}
\end{columns}
\end{frame}

\begin{frame}{2B: STFT - Channel 32}
\fullfig{plot_18.png}
\end{frame}

\begin{frame}{2B: Discussion - Channel 14 vs 32}
\small
\textbf{Signal characteristics (STFT):}
\begin{itemize}
\item The LFM echo is spread equally across $\pm15\,$kHz, starting around $20\,$ms and fading slowly.
\item A weak interfering tone near $15\,$kHz appears around $130\,$ms.
\item The $-17\,$kHz interferer from Channel 14 is barely visible here.
\end{itemize}

\vspace{0.5em}
\textbf{Key difference between channels:}
\begin{itemize}
\item Channel 32 has much less interference from narrowband tones.
\item The interferers that dominate Channel 14 are nearly absent in Channel 32.
\item This spatial interference originates possibly from faulty electronics or near-field source rather than a far-field acoustic source which would appear similarly across all channels.
\end{itemize}
\end{frame}

% =============================================================================
\section*{Summary}
% =============================================================================

\begin{frame}{Summary}
\begin{itemize}
\item \textbf{Ex 1A:} Welch and multitaper reduce variance by $>10\times$ vs the periodogram, with lower passband bias. Estimated bandwidth is $32.1\,$kHz vss. ideal $30\,$kHz; spectral leakage and a $-17\,$kHz interferer bias the estimate.
\item \textbf{Ex 1B:} Variance reduction merges closely spaced tones. High resolution periodogram or windowed periodogram is preferred to resolve narrowband signals.
\item \textbf{Ex 2A:} STFT reveals time-varying events hidden by global PSD. Shorter segments improve time resolution at the cost of frequency resolution.
\item \textbf{Ex 2B:} Channel 32 shows a cleaner LFM echo with minimal interference, suggesting the tones in Channel 14 originate from near-field, likely not far-field sources.
\end{itemize}
\end{frame}

\end{document}
